\section{\texorpdfstring{\code{driver::can::Frame}}{driver::can::Frame}}
\label{sec:framecls}

The register map in \kapref{ch:can} says exactly what the hardware can send and receive, and it is
not much: an identifier, a length and up to eight data bytes. The class mirrors that:

\begin{cppcode}
namespace driver::can
{
/**
 * @brief A single CAN frame: identifier, data length code and payload.
 */
struct Frame
{
    /** Maximum number of payload bytes in a CAN frame. */
    static constexpr std::uint8_t MaxDataLen{8U};

    std::uint16_t id{};               /**< CAN identifier, 11 bits (0x000 - 0x7FF). */
    std::uint8_t  dlc{};              /**< Data length code, 0 - 8. */
    std::uint8_t  data[MaxDataLen]{}; /**< Payload. */
};

/**
 * @brief Check whether a frame is valid, i.e. whether the hardware can send it.
 */
[[nodiscard]] constexpr bool isValid(const Frame& frame) noexcept
{
    return (frame.id <= 0x7FFU) && (frame.dlc <= Frame::MaxDataLen);
}

} // namespace driver::can
\end{cppcode}

That is the whole file. No \code{.cpp}.

\subsection{Why an aggregate}

\code{Frame} deliberately has no constructor, destructor or member function. The reasons, in turn:

\textbf{It is data, not behaviour.} A frame is something that gets moved: from the application to
the driver, from the driver into four registers, from the registers out onto the bus.

\textbf{Aggregate initialisation is readable.} With the aggregate,
\code{Frame frame{0x123U, 2U, {0xAAU, 0xBBU}}} works. Add a single user-declared constructor and
that line stops compiling. The tests you meet from the next chapter onwards are full of such
lines.

\textbf{Default member initialisation makes \code{Frame f{}} safe.} The \code{{}} after every
member means that a default-constructed frame has zeroed fields. Without them, \code{receive()}
could have handed back junk in the bytes it did not fill.

\subsection{\texorpdfstring{\code{static constexpr}}{static constexpr}, not \texorpdfstring{\code{\#define}}{define}}

\code{MaxDataLen} has a \textbf{type}, a \textbf{scope} and exists in the \textbf{debugger}. A
\code{#define} has none of that, and on top of that applies from the line it stands on to the end
of the file, straight through every header included afterwards.

It is also used as the array size, which guarantees that the size of the array and the bounds
check in \code{isValid()} are the same number.

\subsection{Why \texorpdfstring{\code{isValid()}}{isValid()} is a free function}

Two reasons. It keeps the aggregate intact --- a member function is the first step towards a class
with invariants, and the next step is a constructor. And it belongs to the \emph{driver}, not to
the data: the rule comes from this hardware's register widths, not from CAN as a concept.

\begin{aside}
\note[Rule] The hardware validates nothing. The driver \textbf{must} therefore validate --- and
\textbf{only} the driver. The temptation to put one more check in the application "just to be
safe" is to be resisted: two copies of the same rule are two things that have to be kept in step,
and the day \code{MaxDataLen} changes, only one of them will change with it.
\end{aside}
